% !TeX root = ../../Book1.tex
% 纲要标题：### 18.2 迹与范数
% 纲要位置：计划纲要.md，第 623 行。

\section{迹与范数}
\label{b1:sec:1802}

% 纲要内容（仅作写作计划，尚非教材正文）：
%
% * 乘法线性算子的迹与行列式
% * 有限扩张中的域迹与域范数
% * 共轭元公式与不可分重数
%

有限扩张 \(L/K\) 中的元素可以通过“乘以这个元素”的线性变换来观察。这个办法不需要扩张可分，也不需要预先列出共轭根。先用矩阵定义迹与范数，再说明共轭公式成立时根应怎样计数。

\subsection{乘法线性算子的迹与行列式}
\label{b1:subsec:1802:01}

设 \([L:K]=n\)，\(a\in L\)。乘法映射
\[
m_a:L\longrightarrow L,\qquad x\longmapsto ax
\]
是 \(K\) 线性的。选定 \(K\) 基以后，它由一个 \(n\) 阶矩阵表示；换基使矩阵相似，所以矩阵迹、行列式和特征多项式与基的选择无关。矩阵迹在相似变换下不变可由 \(\operatorname{tr}(AB)=\operatorname{tr}(BA)\) 得到，后一个等式只须交换有限求和的指标；行列式的不变性来自乘法公式。

\begin{definition}\label{b1:def:field-trace-norm}
设 \(L/K\) 为有限域扩张，\(a\in L\)。令 \(m_a:L\rightarrow L\) 为 \(K\) 线性映射 \(x\mapsto ax\)。把
\[
\operatorname{Tr}_{L/K}(a)=\operatorname{tr}(m_a),\qquad
\operatorname{N}_{L/K}(a)=\det(m_a)
\]
分别称为元素 \(a\) 关于扩张 \(L/K\) 的\term[yu-ji]{域迹}[field trace]与\term[yu-fanshu]{域范数}[field norm]。映射 \(m_a\) 的特征多项式称为 \(a\) 关于 \(L/K\) 的乘法特征多项式，记为 \(\chi_{a,L/K}(t)=\det(tI-m_a)\in K[t]\)。
\end{definition}
\symidx{\(\operatorname{Tr}_{L/K}\)：有限扩张的域迹}
\symidx{\(\operatorname{N}_{L/K}\)：有限扩张的域范数}
\symidx{\(\chi_{a,L/K}(t)\)：乘以 \(a\) 的特征多项式}
这些量的定义都包含整个扩张 \(L/K\)。它们不是只依赖元素最小多项式的绝对数值。

\subsection{有限扩张中的域迹与域范数}
\label{b1:subsec:1802:02}

\begin{proposition}\label{b1:prop:trace-norm-basic}
设 \(L/K\) 为 \(n\) 次有限域扩张。域迹 \(\operatorname{Tr}_{L/K}:L\rightarrow K\) 是 \(K\) 线性映射。对任意 \(a,b\in L\)，域范数满足
\[
\operatorname{N}_{L/K}(ab)=\operatorname{N}_{L/K}(a)\operatorname{N}_{L/K}(b),
\qquad
\operatorname{N}_{L/K}(a)=0\Longleftrightarrow a=0.
\]
若 \(c\in K\)，则 \(\operatorname{Tr}_{L/K}(c)=nc\)、\(\operatorname{N}_{L/K}(c)=c^n\)。
\end{proposition}
\begin{proof}
等式 \(m_{a+b}=m_a+m_b\)、\(m_{ca}=c m_a\) 和迹的线性给出第一项。又有 \(m_{ab}=m_am_b\)，行列式的乘法公式给出范数的乘法性。若 \(a\ne0\)，映射 \(m_a\) 的逆为 \(m_{a^{-1}}\)，所以其行列式非零；若 \(a=0\)，零矩阵的行列式为零。基域元素对应标量矩阵 \(cI_n\)，给出最后两式。
\end{proof}
设 \(a\) 在 \(K\) 上的首一最小多项式为
\[
f(t)=t^d+c_{d-1}t^{d-1}+\cdots+c_0,
\qquad r=[L:K(a)].
\]
令 \(E=K(a)\)。先把乘法算子限制为 \(m_a|_E:E\rightarrow E\)。在 \(E/K\) 的基 \(1,a,\ldots,a^{d-1}\) 下，这个限制算子将前 \(d-1\) 个基向量依次移向后一个，将末项送到 \(-c_0-\cdots-c_{d-1}a^{d-1}\)，矩阵因而是 \(f\) 的友矩阵 \(C_f\)。沿 \(tI-C_f\) 的末列展开并整理，得到 \(\det(tI-C_f)=t^d+c_{d-1}t^{d-1}+\cdots+c_0=f(t)\)。

\ProseParagraphBreak

再取 \(L/E\) 的一组基 \(v_1,\ldots,v_r\)。由塔式公式的基构造，
\[
v_1,av_1,\ldots,a^{d-1}v_1;\quad\ldots;\quad
v_r,av_r,\ldots,a^{d-1}v_r
\]
是 \(L/K\) 的一组基。每个子空间 \(Ev_j\) 都被 \(m_a:L\rightarrow L\) 保持，并且在所列基下具有同一个矩阵 \(C_f\)。因此，整个乘法矩阵为 \(\operatorname{diag}(C_f,\ldots,C_f)\)，共有 \(r\) 个块，从而
\begin{equation}\label{b1:eq:trace-norm-minimal-polynomial}
\chi_{a,L/K}(t)=f(t)^r,\qquad
\operatorname{Tr}_{L/K}(a)=-r c_{d-1},\qquad
\operatorname{N}_{L/K}(a)=(-1)^{dr}c_0^r.
\end{equation}
例如取非平方的 \(D\in\mathbb Q\)，令 \(L=\mathbb Q(\sqrt D)\)，并取 \(a,b\in\mathbb Q\)。乘以 \(a+b\sqrt D\) 的矩阵为
\[
\begin{pmatrix}a&bD\\ b&a\end{pmatrix},
\]
所以迹为 \(2a\)，范数为 \(a^2-Db^2\)。当 \(D<0\) 时，这个范数与复共轭有关，但一般域范数并不是分析学意义的非负长度。

\subsection{共轭元公式与不可分重数}
\label{b1:subsec:1802:03}

\begin{theorem}\label{b1:thm:trace-norm-embeddings}
设 \(L/K\) 为有限域扩张，\(\Omega\) 为包含 \(L\) 的 \(K\) 的代数闭包，\(\sigma_1,\ldots,\sigma_s\) 是全部不同的 \(K\) 嵌入 \(L\rightarrow\Omega\)。记 \(e=[L:K]_{\mathrm i}=[L:K]/s\) 为不可分次数。则对每个 \(a\in L\) 都有
\begin{align*}
\chi_{a,L/K}(t)&=\prod_{i=1}^s\bigl(t-\sigma_i(a)\bigr)^e,\\
\operatorname{Tr}_{L/K}(a)&=e\sum_{i=1}^s\sigma_i(a),&
\operatorname{N}_{L/K}(a)&=\prod_{i=1}^s\sigma_i(a)^e.
\end{align*}
特别地，扩张可分时 \(e=1\)，迹是全部嵌入像的和，范数是其积。
\end{theorem}
\begin{proof}
令 \(E=K(a)\)、\(f=m_{a,K}\)、\(d=[E:K]\)。最小多项式 \(f\) 有 \(d_s=[E:K]_{\mathrm s}\) 个不同根；\cref{b1:prop:inseparable-polynomial-shape}说明正特征下各根重数相同，特征零下各根重数为一，所以统一为 \(d/d_s\)。每个 \(K\) 嵌入 \(E\rightarrow\Omega\) 都能延拓到 \(L\)，且延拓的个数恰为 \(r_s=[L:E]_{\mathrm s}\)，这是\cref{b1:thm:separable-degree-tower}证明中的嵌入计数。因此 \(s=d_s r_s\)，而 \(\sigma_i(a)\) 中每个不同根出现 \(r_s\) 次。令 \(r=[L:E]\)，则
\[
e r_s=\frac{dr}{d_s r_s}r_s=\frac d{d_s}r.
\]
右端恰是 \(f(t)^r\) 中每个不同根的重数。由\eqref{b1:eq:trace-norm-minimal-polynomial}，这就证明特征多项式公式。比较 \(t^{[L:K]-1}\) 的系数与常数项，分别得到迹和范数的公式。
\end{proof}
公式中的重复有两个来源，即不同嵌入在指定元素上取同一值，以及每个嵌入像还要按不可分次数 \(e\) 计重。即使扩张可分，第一种重复也可能出现。取 \(L=\mathbb Q(\sqrt2,\sqrt3)\)、\(K=\mathbb Q\)、\(a=\sqrt2\)。分别为 \(\sqrt2\) 和 \(\sqrt3\) 选择正负像，得到四个不同的 \(K\) 嵌入；它们在 \(a\) 上的取值为 \(\sqrt2,\sqrt2,-\sqrt2,-\sqrt2\)。此时 \(e=1\)，但 \(\chi_{a,L/K}(t)=(t^2-2)^2\)。因此，求迹或范数时仍须保留不同嵌入带来的重复取值。

\ProseParagraphBreak

另一种重复来自不可分次数。设 \(p\) 为素数，取在 \(\mathbb F_p\) 上超越的元素 \(u\)，并令 \(L=\mathbb F_p(u)\)、\(K=\mathbb F_p(u^p)\)。由\eqref{b1:eq:rational-power-degree}知 \([L:K]=p\)；每个 \(x\in L\) 都满足 \(x^p\in K\)，故这是纯不可分扩张，只有一个 \(K\) 嵌入，而 \(e=p\)。于是
\[
\chi_{u,L/K}(t)=(t-u)^p=t^p-u^p,\qquad
\operatorname{Tr}_{L/K}(u)=pu=0,\qquad
\operatorname{N}_{L/K}(u)=u^p.
\]
事实上，这个扩张中每个元素的迹都是零。唯一嵌入的像若不按 \(e\) 计重，就无法得到这些结果。
